% Paragraph for the paper: what Lean's kernel checks in the kernel-only formalization.
% Numbers from logs/run2/ of the Lean artifact (S13Kernel). Uses the bibliography keys of
% lean_bib.tex (Lean4, Mathlib, LeanRepo) from the native_decide project.

\paragraph{Kernel-checked computation.}
The computation is also checked without native code. An untrusted program produces a certificate: a
list $H$ of $99{,}736$ states, the start state among them, and for each state and each comparator with
$\mathrm{lo}<\mathrm{hi}$ whose next state is neither unchanged nor pruned by the potential $\Phi$, a
permutation of the channels and the index of the state of $H$ that the relabeled next state equals.
Of the $99{,}736\cdot78=7{,}779{,}408$ next states, $275{,}115$ are unchanged, $6{,}396{,}523$ have
$\Phi>2^8$ and $1{,}107{,}770$ are relabelings of states of $H$. A checking function written for Lean's
kernel recomputes every next state and its potential, verifies each permutation and each relabeled
state entry by entry, and confirms that no state of $H$ is a goal; Lean proves that its success gives
the hypotheses of the closure argument, hence $P(2,13)\ge9$. Lean's kernel evaluates the check itself
(\texttt{decide +kernel}) in $1{,}102$ theorems of about $90$ states each. Every theorem of the proof,
including \texttt{main\_theorem}, therefore depends only on the standard axioms:
\begin{alltt}\small
'S13K.main_theorem' depends on axioms: [propext, Classical.choice, Quot.sound]
\end{alltt}
and the proof trusts Lean's kernel together with the hypothesis $S(11)\ge35$; the certificate and the
program that produced it are not trusted. The statement is the same Lean proposition as above. On a
$384$-thread x86-64 server, building the proof with $120$ parallel jobs took $2$~min~$56$~s of wall
clock time and $4.96$ hours of processor time, with at most $4.2$~GiB of memory per job (Lean~4
\texttt{v4.35.0-rc3}). Each of seven corrupted certificates (a wrong index, a wrong permutation, a missing
entry, a changed state) is rejected.
